\chapter{Differential Diagnosis}
\section*{What Is This Test?} Differential diagnosis is the structured process of comparing possible causes of a symptom or finding. It is not a single laboratory test.
\section*{Alternative Names} Diagnostic workup; clinical differential; diagnostic reasoning.
\section*{Principle} History, examination, prevalence, risk, and test characteristics are combined to rank competing diagnoses and select efficient testing.
\section*{Why This Test Is Ordered} The concept applies whenever symptoms have multiple possible causes or an initial result is nonspecific.
\section*{Primary vs Secondary Test} Clinical history and examination establish the initial differential; targeted laboratory, imaging, pathology, and physiologic tests refine it.
\section*{How This Test Is Performed} The clinician identifies urgent threats, gathers focused information, orders tests, updates probabilities, and reassesses the patient as results arrive.
\section*{What Preparation Is Needed} Preparation depends on the selected test. Provide complete symptom, exposure, medicine, and family history.
\section*{Understanding Results} No test result should be interpreted without pretest probability, limitations, false positives, false negatives, and the patient's clinical course.
\section*{Follow-up Tests} Follow-up is diagnosis-specific and may include repeat examination, confirmatory testing, referral, or observation.
\section*{References} \begin{enumerate}\item MedlinePlus. Differential Diagnosis.\item Agency for Healthcare Research and Quality. Diagnostic safety resources.\end{enumerate}
\clearpage\section*{Understanding Results and Follow-up}
The laboratory report should identify the specimen, method, units, reference interval or decision limit, and whether the result is positive, negative, abnormal, or inconclusive. Reference intervals vary with age, sex, pregnancy, specimen type, assay, and laboratory; a result outside a range is not automatically a diagnosis, and a result within a range does not always exclude disease. Discuss unexpected results with the ordering clinician, who can decide whether repeat or confirmatory testing, treatment, or referral is appropriate. Do not change medicines or delay urgent care based on an isolated result.
\section*{Regulatory Status and Authorization Date}This chapter describes a test category, not a specific commercial product. FDA, CDSCO, EU IVDR, and MHRA approval or registration dates therefore cannot be assigned without the assay manufacturer, product name, instrument, and intended use. For laboratory-developed tests, an individual agency approval date may not exist. See the Preface for the interpretation of regulatory status.
\clearpage
